\hypertarget{citations}{%
\section{Citations}\label{citations}}

Hello world，世界你好。 \citeauthor{coser1956}
\citetext{\citeyear{coser1956}; \citeyear{coser1957}; \citeyear{coser1967}}
认为冲突能够在社会关系和社会结构中发挥积极作用，例如不同阶层、民族和党派之间的冲突，能够起到建立和维护自身群体身份的作用，把其他群体排除在外。此外，敌对状态往往也通过冲突的形式表现出来，换句话说，冲突是敌对的一种外在表现，以维持各个主体之间的利益平衡关系。从这个角度来看，以实现具体目标或利益的战略行为为特征的``理性主义''或``现实主义''的冲突，更有可能产生建设性的结果
\citep{coser1956, deutsch1969}。然而，以高度紧张、压抑和敌意为特征的``非理性主义''或``超现实主义''社会冲突，可能会导致破坏性的结果
\citep{deutsch1969}。在灵活性的社会结构中，冲突可以被允许存在，甚至被制度化，用来化解社会不满。相反僵化的社会结构缺乏减少或化解冲突的机制，导致冲突可能会被抑制，累积到一定程度，将使社会系统崩溃
\citep{sidaway2005}。

综合这两类观点，我们可以认为，冲突既是变化的结果，也是变化的来源。制度与冲突之间的联系是在一定程度上的冲突，因此
\citeauthor{ho2006e}
\citetext{\citeyear{ho2006e}; \citeyear{ho2013}; \citeyear{ho2014}; \citeyear{ho2016a}}
认为，可以通过各个主体对冲突的感知来测量制度的可信度，即制度可信度理论。制度可信度理论认为，超过一定程度的冲突后，原有的制度将无法解决冲突，而新的制度安排将会出现，因此从这一观点出发，制度安排下冲突越多，制度的可信度就越低。使用冲突的程度来作为制度可信度的测量指标，其重点是关注冲突本身。

\hypertarget{table}{%
\section{Table}\label{table}}

\begin{longtable}[]{@{}cccc@{}}
\caption{预期成果列表}\tabularnewline
\toprule
序号 & 成果形式 & 完成时间 & 使用去向 \\ \addlinespace
\midrule
\endfirsthead
\toprule
序号 & 成果形式 & 完成时间 & 使用去向 \\ \addlinespace
\midrule
\endhead
1 & 学术论文一 & 2021年9月 & 发表第一篇SSCI期刊论文 \\ \addlinespace
3 & 学术座谈会 & 2022年3月 & 媒体顾问文章 \\ \addlinespace
4 & 学术论文二 & 2022年6月 & 发表第二篇CSSCI期刊论文 \\ \addlinespace
5 & 学术论文三 & 2022年9月 & 发表第三篇SSCI期刊论文 \\ \addlinespace
6 & 书籍 & 2023年10月 & 在国际著名出版社公开出版 \\ \addlinespace
7 & 政策咨询报告 & 2023年10月 &
供中央和部分地方政府作为政策参考 \\ \addlinespace
8 & 国际会议 & 2023年11月 & 会议论文 \\ \addlinespace
\bottomrule
\end{longtable}

\hypertarget{comments}{%
\section{Comments}\label{comments}}

\hypertarget{images}{%
\section{Images}\label{images}}

\begin{figure}
\centering
\includegraphics[width=0.7\textwidth,height=0.5\textheight]{images/dog.jpg}
\caption{This is a dog}
\end{figure}
